Every element of $r(\mathfrak{a})$ vanishes on $V(\mathfrak{a})$. Conversely, suppose $f\notin r(\mathfrak{a})$ and choose a prime $\mathfrak{p}\supseteq\mathfrak{a}$ with $f\notin\mathfrak{p}$. The ring $C=(k[t_1,\ldots,t_n]/\mathfrak{p})_f$ is nonzero and finitely generated over $k$. Choose a maximal ideal $\mathfrak{m}\subseteq C$. By the weak Nullstellensatz, $C/\mathfrak{m}\cong k$. The images of the $t_i$ define a point of $V(\mathfrak{a})$, while the image of $f$ is nonzero because it is a unit in $C$. Thus $f$ does not vanish on all of $V(\mathfrak{a})$.
